\documentclass[../DoAn.tex]{subfiles}
\begin{document}

\begin{center}
    \Large{\textbf{ABSTRACT}}\\
\end{center}
\vspace{1cm}
This graduation thesis focuses on building a chatbot system to support sales consultation, product comparison, and price reference in the furniture domain. In practice, small and medium-sized furniture stores need to advise customers effectively, but they often lack the technical resources to develop their own AI-based systems. In addition, furniture products usually contain many important attributes, such as dimensions, materials, styles, and prices. Customers may also change their requirements, budgets, or reject previous suggestions during a conversation. Existing solutions, including rule-based chatbots, RAG-based chatbots, and commercial AI chatbots, still have certain limitations. They may struggle with flexible multi-turn conversations, require technical effort to build and maintain a knowledge base, or incur high operating costs.

In this thesis, I choose to develop a multi-tenant RAG chatbot system, in which each store has its own data, chatbot, and configuration while sharing the same underlying platform. This approach is suitable for furniture sales consultation because it allows the system to retrieve knowledge from each store’s product data and, at the same time, support multiple stores in a scalable manner.

The proposed solution consists of a backend for managing tenants, chatbots, users, conversations, knowledge bases, and purchase requests; an AI service for handling dialogue, retrieving relevant knowledge, and generating consultation responses; an administration interface for stores and the system; and a chat interface for end users. The system is designed to support three main modes: store-specific consultation, general consultation and product comparison, and market price reference. The main contribution of this thesis is the design and implementation of a multi-tenant RAG chatbot system for furniture sales consultation, with the ability to support multi-turn conversations, recommend and compare products based on data, and provide a foundation for evaluating retrieval quality, response quality, and consultation effectiveness.

\end{document}